\documentclass{article}

\usepackage{amsthm}
\usepackage{graphicx}
% \usepackage{amssymb}

\newtheorem{proposition}{Proposition}

\title{MATH290 Problem 1}
\author{Caleb Hessing}
\date{September 6, 2026}

\begin{document}
\maketitle 

A famous pirate captain has hidden his treasure in three chests. One chest contains
gold coins, one chest contains counterfeit coins, and one chest contains a mixture of
gold and counterfeit coins. Each chest is labeled to indicate its contents: either G
for gold, C for counterfeit, or B for both. Because this pirate captain is very clever
and dishonest, none of the labels match the actual contents of the chest. You cannot
visually distinguish a gold coin from a counterfeit one, but you can perform a lengthy
chemical test to tell the two apart. What is the minimum number of coins you need to
test to determine the contents of all three chests, and what strategy should you use?

\begin{proposition}
    Only one coin is necessary to determine the contents of all three chests.
\end{proposition}

\begin{proof}
Hello there.

\end{proof}

\end{document}